\begin{abstract}
Ransomware and credential-stealing malware leave their most useful evidence in RAM: injected code, the
ciphers they load, live keys and session tokens. Existing tools each answer one part of the question, and
large language models (LLMs) that could write the incident report are both prone to unsupported claims and
exposed to prompt injection through the evidence itself. We present TriageGuard, a pipeline that takes a
Windows or Linux memory image, extracts processes and injected code with Volatility~3, identifies ciphers with
a Dempster-Shafer confidence, tests every RSA key against NIST SP~800-57, recovers live AES, ChaCha20 and RSA
private keys, decodes tokens, and fuses the evidence into a risk score whose every input is cited. An LLM
writes the report from the numbered findings only; four defences (sanitising, spotlighting, a citation check
and a rule-based cross-check) gate what reaches the reader. On 14 public and self-built images the pipeline
separates a clean baseline (risk 28) from compromised hosts (67 to 97), rebuilds four RSA-2048 private keys
from a WannaCry image, and passes all 12 ground-truth checks on a VM with planted keys. A benchmark of 50
processes identifies the cipher in use 45/45 times. Against 15 injection attacks, including adaptive ones,
2/90 succeed undefended and 0/90 defended, and none reaches a reader unflagged. All 14 LLM reports pass
validation with 90 to 100\% of claims cited and the stated risk equal to the computed one.
\end{abstract}

\begin{IEEEkeywords}
memory forensics, cryptographic key recovery, RSA weak keys, Dempster-Shafer evidence fusion, prompt injection, large language models
\end{IEEEkeywords}
